\section{Introduction}\label{sec:introduction}
Obfuscation is common in investor-directed fund communications.
Even when pursuing simple strategies, fund managers describe them in markedly different and often unnecessarily complex ways. For example, as \citet*{dehaan2021obfuscation} document, two funds that both track the S\&P 500 describe essentially identical objectives with strikingly different clarity: Schwab states in a single sentence, “The fund's goal is to track the total return of the S\&P 500 Index,” while Deutsche describes it in 60 words, embedding the same objective in technical language.\footnote{Deutsche's S\&P 500 index fund is described as “The fund seeks to provide investment results that, before expenses, correspond to the total return of common stocks publicly traded in the United States, as represented by the Standard \& Poor's 500 Composite Stock Price Index (S\&P 500 Index). The fund invests for capital appreciation, not income; any dividend and interest income is incidental to the pursuit of its objective.”} Similar patterns arise in fund names: labels such as “Dynamic Opportunity Fund,” “Flexible Growth Strategy,” or “Opportunistic Equity Fund” convey little concrete meaning to unprofessional retail investors at first glance, yet in practice often correspond to relatively simple large-cap portfolios that closely resemble the S\&P 500. This pattern is puzzling, as standard signaling theory suggests that managers, seeking to convey skill and attract capital, should prefer greater transparency rather than opacity.
Why would managers make information harder for their own investors to interpret, and what are the incentives and market consequences when fund opacity becomes a manager's strategic choice?

To address these questions, we study an equilibrium model of financial market trading with delegated asset management. A fund manager, privately informed about her skill, raises capital from investors through obfuscated communications and trades a risky asset in the financial market \`a la \citet{kyle1985continuous}. We define skill as the manager-specific ability to uncover more information about asset payoffs than the market: the more skilled the manager, the larger the mispricing she can identify and the greater the trading profits she can potentially earn. To attract fund flows, the manager privately communicates information about her risky position (portfolio) to investors. Investors, in turn, use this fund information to infer the manager's skill and decide on capital allocation to the fund. Importantly, the manager strategically chooses the readability of this communication: by obfuscating information, she makes it more difficult for investors to back out her underlying skill.\footnote{Following \citet*{ellison2009search}, the term ``obfuscation" refers to general practices that reduce the readability of information and distort communications. Accordingly, ``transparency" and ``opacity" refer to the degre of interpretability and precision of information with which investors can assess the fund’s positions, rather than mere disclosure.}

The model yields two stable, self-fulfilling equilibria regarding the manager's trading strategy. One of the equilibria is similar to that of \citet{kyle1985continuous}, where the manager employs the profit-maximizing trading strategy. In the other, which we call the “flow-driven” equilibrium, the manager trades the risky asset more aggressively to appear skilled and to attract large fund flows. This strategy increases trading volume, price volatility, and market liquidity but undermines the fund expected performance by revealing too much private information to the market maker.
Intuitively, the manager can boost investors' perception of her skill by placing a large order, either buy or sell, because such an action makes her appear like a skilled manager who has identified an asset with significant mispricing.
%\footnote{In equilibrium, investors accurately anticipate the manager's strategy, so their learning is not manipulated; nevertheless, the manager may take excessive risk following investors' belief that the manager trades more aggressively than the profit-maximizing level, as doing otherwise makes investors (incorrectly) infer a less skilled manager and reduce fund flows.} 
Although the expected performance deteriorates, it triggers large fund flows, and the manager earns a large fee income. Key to our multiple equilibria is the manager's conflicting motives for revealing skill: on one hand, she wishes to hide her skill from the market maker to earn trading profits, but on the other hand, she wants to show it off to investors to establish high reputation and attract large fund flows. 

The relative strength of these competing motives is governed by the degree of fund opacity. When opacity is high, investors can extract little information from fund communications and do not rely on them to infer skill. Thus, flows are insensitive to the manager's trading strategy, the showing-off motive is suppressed, and only the Kyle-like equilibrium survives. When opacity is low, investors rely heavily on fund information, and only the flow-driven equilibrium exists. At intermediate levels of opacity, however, both equilibria coexist, and which one prevails is determined by investors' beliefs. If investors expect aggressive, flow-driven trading, they treat fund information as a reliable signal of skill, making flows highly responsive to reported positions. It strengthens the showing-off motive and induces the manager to trade aggressively, confirming the investors' beliefs. The symmetric argument sustains the Kyle-like equilibrium when investors expect profit-maximizing behavior. A mere shift in beliefs, unaccompanied by any change in fundamentals, can therefore trigger a transition from the Kyle-like to the flow-driven equilibrium, where price volatility is high and trading performance is poor. We interpret this susceptibility to non-fundamental belief shifts as \textit{fragility}.\footnote{As in \citet{greenwood2011stock}, we define a market as fragile if it is susceptible to non-fundamental demand shifts, such as those caused by changes in investor beliefs.} 

%The equilibrium multiplicity in trading styles has implications for the fund-management industry. For example, the flow-driven equilibrium explains ovservations by \citet{philippon2012wages} and \citet{celerier2019returns}, who document increasing returns to talent and their significant contribution to income inequality in the finance industry. In the flow-driven equilibrium, managers' fee income exhibits greater inequality among high- and low-skilled managers compared to those in the Kyle-like equilibrium, as they trade aggressively and reveal their skills more than they do in the Kyle-like equilibrium. This allows investors to better assess each manager's skill and, therefore, to allocate capital more accurately according to managers' skills. Moreover, multiple equilibria offer a novel explanation for the puzzling relationship (or the lack thereof) between skill and persistent performance. In our model, skill does not translate into persistent performance, as high-skilled managers in the flow-driven equilibrium may underperform low-skilled ones following the Kyle-like equilibrium, consistent with the evidence documented in the empirical literature (e.g., \citealp*{busse2010performance}; \citealp{berk2004mutual}).

Building on these results, we endogenize fund opacity and study which types of managers contribute to financial fragility. By obfuscating fund information, a manager weakens investors' responsiveness to reported positions and curbs her showing-off incentive at the portfolio-decision stage. Although, when deciding on her risky position, she takes investors' belief-updating rule as given and is tempted to over-trade to attract flows, she can \textit{ex-ante} influence that rule through obfuscation. Opacity thus serves as a commitment device: by tying her own hands, the manager mitigates inefficient risk-taking, raises the profit margin of her risky position, and improves the fund's expected performance, thereby benefiting both herself and investors.

The value of this commitment is greater for higher-skilled managers, as they benefit more from improved profit margins. Optimal opacity therefore increases with managers' skill. This result is consistent with the ``brain drain'' documented by \citet{kostovetsky2017brain}, whereby top mutual fund managers migrate to more opaque hedge funds and deliver superior performance. Crucially, the relationship between opacity and fragility is non-monotone. While high opacity yields a unique Kyle-like equilibrium and low opacity yields a unique flow-driven equilibrium, it is at intermediate levels of opacity where multiple equilibria, and thus fragility, arise. Because intermediate-skilled managers optimally choose intermediate opacity, they expose financial markets to belief-driven instability. Such behavior is consistent with momentum crashes and sudden reversals observed during financial crises \citep{brunnermeier2009deciphering}, as well as the abrupt strategy unwinding among quantitative funds documented by \citet{khandani2011what}. Importantly, fragility in our model is driven entirely by shifts in market beliefs and can arise without fundamental shocks, in line with the argument that large historical price swings are rarely accompanied by proportionally large changes in fundamentals \citep{gennotte1999market}. A direct policy implication follows: transparency mandates may push managers toward lower opacity and move high-skilled managers into the fragile intermediate region, thereby \textit{increasing} rather than decreasing systemic fragility. Transparency, in this sense, backfires.

%% Literature Review %%
Our paper contributes to the understanding of opacity, complexity, and obfuscation in financial markets. Empirically, \citet*{edelen2012disclosure}, \citet*{badoera2020icansee}, and \citet*{dehaan2021obfuscation} document that mutual funds strategically obfuscate disclosures of their strategies and fee sturcutures, and \citet{celerier2019returns} show that structured financial products are deliberately made complex. Building on such evidence, \citet{carlin2009strategic}, \citet{carlin2010obfuscation}, and \citet{sato2014opacity} develop theoretical models in which complexity enables financial institutions to extract fees or preserve informational advantages. In these models, however, opacity is typically treated as technologically given or engineered externally, while we analyze its implications when fund managers strategically obfuscate. Moreover, opacity in our framework improves both managers’ and investors’ \textit{ex-ante} utility. Unlike the prior literature, managers obfuscate not to exploit less-informed investors, but because opacity serves as a commitment device that mitigates managers' showing-off motive in financial markets.

Our paper is also related to the literature on fragility and trading frenzies in financial markets. \citet{cespa2015beauty} is most closely related, as they also find multiple equilibria arising from short-term investor horizons and persistent liquidity trading. In \citet*{froot1992herd}, strategic complementarities among short-term traders are also important for the emergence of multiple equilibria. Unlike these studies, our model does not feature short horizons; instead, we obtain fragility from managers conserns about fund flows. In \citet{cespa2014illiquidity}, learning from other assets amplifies shocks and causes liquidity crashes. Although the mechanism is different, our model can also result in liquidity crashes when the economy switches between equilibria due to non-fundamental shocks. Also, several studeits, such as \citet{veldkamp2006media} and \citet*{goldstein2013trading}, features trading frenzies based on strategic complementarities in information markets or a feedback effect from financial markets to a firm's real investment. By contrast, in our flow-driven equilibrium, a fund manager trades aggressively out of fear of being perceived as less skilled by the recruiter who believes that she would trade aggressively. 

Lastly, our work contributes to the broader understanding of the role of asymmetric information about skills in financial markets. \citet{Huberman1993ontheincentives}, \citet{huddart1999reputation}, \citet{di2015fake}, \citet*{malliaris2021reputation}, and \citet*{bijlsma2018competition} derive implications for the risk-taking behavior of agents when ability is private information but is conveyed to investors through signals, such as realized performance. In the context of the principal-agent proglem, \citet{prat2005wrong} shows that transparency in such communications can be detrimental, as it distorts agents' behavior. We also link trader skill to private information and asset markets; however, unlike these papers, our focus is on deriving pricing implications and managers' incentive to obfuscate fund information. \citet{guerrieri2012fund} show that career concerns amplify price volatility, similar to our findings. In \citet{dasgupta2006financial}, career concerns lead to churning (trading without private information), and in \citet{dasgupta2008information}, they lead to herding (ignoring private information and trading like everyone else). These three papers are closely related to ours in that they study career concerns and endogenous asset prices, but they do not analyze fragility and endogenous fund opacity, the main focus of our paper. \citet{gervais2020transparency} consider managers’ transparency choices under exogenous transparency costs and obtain sorting across skill types. By contrast, we derive an endogenous cost of transparency arising from flow-driven trading distortions. As a result, transparency in our model does not merely separate types but alters trading behavior and can generate belief-driven multiple equilibria and market fragility.

The rest of the paper proceeds as follows. Section \ref{sec:Model} presents a baseline model with exogenous reputation concerns, and Section \ref{sec:eqm} studies the equilibrium. Section \ref{sec:endogenous} endogenizes reputation concerns by introducing a two-period model and presents the implications of our results. Section \ref{sec:skillacquisition} analyzes traders' skill acquisition. The \hyperref[app:appendix]{Online Appendix} contains all proofs for the theoretical results.